\documentclass[12pt,a4paper]{article}
\usepackage[utf8]{inputenc}
\usepackage[T1]{fontenc}
\usepackage[spanish]{babel}
\usepackage{geometry}
\usepackage{graphicx}
\usepackage{enumitem}
\usepackage{booktabs}
\usepackage{array}
\usepackage{xcolor}
\usepackage{hyperref}
\usepackage{fancyhdr}
\usepackage{titlesec}
\usepackage{tcolorbox}

\geometry{margin=2.5cm}
\hypersetup{colorlinks=true, linkcolor=blue!60!black, urlcolor=blue!60!black}

% Colores
\definecolor{maravia}{RGB}{0, 102, 153}
\definecolor{gris}{RGB}{245, 245, 245}
\definecolor{verde}{RGB}{34, 139, 34}
\definecolor{rojo}{RGB}{200, 50, 50}

% Estilos de sección
\titleformat{\section}{\Large\bfseries\color{maravia}}{}{0em}{}[\titlerule]
\titleformat{\subsection}{\large\bfseries\color{maravia!80!black}}{}{0em}{}

% Encabezado
\pagestyle{fancy}
\fancyhf{}
\fancyhead[L]{\small\color{gray}MaravIA Bot --- Manual de Usuario}
\fancyhead[R]{\small\color{gray}\thepage}
\renewcommand{\headrulewidth}{0.4pt}

% Cajas de ejemplo
\newtcolorbox{mensajebot}{colback=gris, colframe=maravia, title=Bot responde:, fonttitle=\bfseries\small}
\newtcolorbox{mensajeusuario}{colback=green!5, colframe=verde!70!black, title=Usuario escribe:, fonttitle=\bfseries\small}
\newtcolorbox{notaimportante}{colback=red!5, colframe=rojo, title=Importante, fonttitle=\bfseries\small}

\begin{document}

% ============================================================
% PORTADA
% ============================================================
\begin{titlepage}
\centering
\vspace*{3cm}
{\Huge\bfseries\color{maravia} MaravIA Bot \par}
\vspace{0.5cm}
{\Large Manual de Usuario \par}
\vspace{1cm}
{\large Registro de Ventas y Compras por WhatsApp \par}
\vspace{2cm}
{\normalsize
Sistema de registro contable conversacional \\
mediante inteligencia artificial \\[1cm]
Plataforma Maravia --- Per\'u
\par}
\vfill
{\small Versi\'on 1.0 --- Marzo 2026 \par}
\end{titlepage}

\tableofcontents
\newpage

% ============================================================
% INTRODUCCIÓN
% ============================================================
\section{Introducci\'on}

MaravIA Bot es un asistente de WhatsApp que permite registrar operaciones contables (ventas y compras) de forma conversacional. El bot gu\'ia al usuario paso a paso para completar todos los datos necesarios y emitir el comprobante electr\'onico.

\subsection{Requisitos}

\begin{itemize}
    \item Un n\'umero de WhatsApp vinculado a la plataforma Maravia
    \item Acceso a la empresa configurada en el sistema
    \item Los datos de la operaci\'on (cliente/proveedor, monto, tipo de documento, etc.)
\end{itemize}

% ============================================================
% FLUJO GENERAL
% ============================================================
\section{Flujo general}

El proceso de registro tiene \textbf{4 etapas} principales:

\begin{enumerate}[label=\textbf{Etapa \arabic*:}, leftmargin=3.5cm]
    \item \textbf{Inicio} --- Indicar si es una venta o una compra
    \item \textbf{Datos del comprobante} --- Proporcionar cliente/proveedor, monto, tipo de documento, moneda y m\'etodo de pago
    \item \textbf{Opciones operativas} --- Elegir sucursal, centro de costo (solo compras) y forma de pago
    \item \textbf{Emisi\'on} --- Confirmar y emitir el comprobante
\end{enumerate}

El bot acepta m\'ultiples formatos de entrada:

\begin{itemize}
    \item \textbf{Texto}: mensajes escritos con los datos de la operaci\'on
    \item \textbf{Audio}: notas de voz que el sistema transcribe autom\'aticamente
    \item \textbf{Documento (PDF)}: se almacena como documento adjunto al registro
    \item \textbf{Imagen}: se almacena como documento adjunto al registro
\end{itemize}

Si el usuario env\'ia un documento o imagen, este se guardar\'a autom\'aticamente como \textbf{documento adjunto} a la operaci\'on y se vincular\'a al comprobante al finalizar.

El bot solo pregunta por datos que faltan. Si el usuario proporciona varios datos en un solo mensaje, el bot los registra todos y solo pregunta por lo pendiente.

% ============================================================
% REGISTRO DE VENTA
% ============================================================
\section{Registro de una venta}

\subsection{Paso 1: Iniciar la conversaci\'on}

\begin{mensajeusuario}
Hola
\end{mensajeusuario}

\begin{mensajebot}
\textit{El bot saluda y muestra dos botones:}
\begin{itemize}[label=\textbullet]
    \item Registrar compra
    \item Registrar venta
\end{itemize}
\end{mensajebot}

El usuario tambi\'en puede escribir directamente: \texttt{"Quiero registrar una venta"} sin necesidad de los botones.

\subsection{Paso 2: Proporcionar los datos}

El bot pedir\'a los datos que falten. Se pueden enviar en cualquier orden y en uno o varios mensajes.

\subsubsection{Datos obligatorios}

\begin{center}
\begin{tabular}{>{\bfseries}l p{8cm}}
\toprule
Campo & Descripci\'on \\
\midrule
Cliente & Nombre o raz\'on social del cliente \\
RUC / DNI & Documento del cliente (8 o 11 d\'igitos) \\
Tipo de documento & Factura, boleta o nota de venta \\
Monto & Monto total de la operaci\'on \\
Moneda & Soles (PEN) o d\'olares (USD) \\
M\'etodo de pago & Contado o cr\'edito \\
\bottomrule
\end{tabular}
\end{center}

\subsubsection{Identificaci\'on del cliente o proveedor}

El bot necesita identificar al cliente (venta) o proveedor (compra) mediante su \textbf{RUC o DNI}. Existen tres escenarios:

\begin{enumerate}
    \item \textbf{El RUC/DNI ya est\'a registrado}: el bot devuelve autom\'aticamente el nombre asociado. No es necesario escribirlo.

    \item \textbf{El RUC/DNI no est\'a registrado}: si env\'ias el n\'umero junto con el nombre de la persona o empresa, el bot lo registrar\'a autom\'aticamente para futuras operaciones.

    \item \textbf{Error tipogr\'afico en el nombre}: si al registrar se cometi\'o un error en el nombre, la correcci\'on debe hacerse desde la \textbf{p\'agina web de Maravia} editando al cliente o proveedor directamente.
\end{enumerate}

\begin{mensajeusuario}
El cliente es Juan P\'erez, DNI 12345678
\end{mensajeusuario}

\begin{mensajebot}
\texttt{Cliente encontrado: Juan P\'erez (12345678).}
\end{mensajebot}

Si no est\'a registrado:

\begin{mensajebot}
\texttt{Cliente registrado: Juan P\'erez (12345678).}
\end{mensajebot}

\begin{notaimportante}
Si el m\'etodo de pago es \textbf{cr\'edito}, el bot preguntar\'a adicionalmente:
\begin{itemize}
    \item \textbf{D\'ias de cr\'edito}: 15, 30, 45, 60 o 90 d\'ias
    \item \textbf{N\'umero de cuotas}: de 1 a 24
\end{itemize}
Estas preguntas \textbf{no aparecen} si el m\'etodo es contado.
\end{notaimportante}

\subsubsection{Ejemplo de conversaci\'on}

\begin{mensajeusuario}
Venta de factura por 1500 soles a Empresa SAC, RUC 20123456789, contado
\end{mensajeusuario}

\begin{mensajebot}
\textit{El bot registra todos los datos de una vez y muestra:}

\medskip
\texttt{Anotado.}

\medskip
\texttt{VENTA} \\
\texttt{Factura} \\
\texttt{Cliente: Empresa SAC} \\
\texttt{RUC: 20123456789} \\
\texttt{Total: S/ 1,500.00} \\
\texttt{Moneda: PEN} \\
\texttt{M\'etodo: Contado}

\medskip
\texttt{Cliente encontrado: Empresa SAC (20123456789).}

\medskip
\texttt{Confirmar todo para continuar?}
\end{mensajebot}

Si faltaran datos, el bot mostrar\'ia:

\begin{mensajebot}
\texttt{Por favor, br\'indame estos datos:} \\
\texttt{1. Moneda: en soles (PEN) o d\'olares (USD)?} \\
\texttt{2. M\'etodo de pago: contado o cr\'edito?}
\end{mensajebot}

\subsection{Paso 3: Confirmar el registro}

Cuando todos los datos est\'an completos, el usuario confirma:

\begin{mensajeusuario}
Confirmo
\end{mensajeusuario}

Palabras aceptadas: \textit{s\'i, confirmo, dale, correcto, listo, ok, de acuerdo, va, perfecto, adelante, acepto, vale, est\'a bien, procede}.

\subsection{Paso 4: Elegir opciones operativas}

Tras confirmar, el bot presenta listas de WhatsApp para elegir:

\begin{enumerate}
    \item \textbf{Sucursal} --- Lista de sucursales de la empresa
    \item \textbf{Forma de pago} --- Transferencia, efectivo, tarjeta, billetera virtual, etc.
\end{enumerate}

\begin{mensajebot}
\texttt{Por favor, br\'indame la sucursal dentro de las siguientes opciones:} \\
\texttt{- Sucursal Principal} \\
\texttt{- Sucursal Norte} \\
\texttt{- Sucursal Sur}
\end{mensajebot}

\begin{mensajeusuario}
Sucursal Principal
\end{mensajeusuario}

Despu\'es de elegir todas las opciones:

\begin{mensajebot}
\texttt{Por favor, br\'indame una confirmaci\'on para continuar}
\end{mensajebot}

\subsection{Paso 5: Confirmaci\'on final y emisi\'on}

\begin{mensajeusuario}
Dale
\end{mensajeusuario}

\begin{mensajebot}
\texttt{VENTA REGISTRADA EN SUNAT!}

\medskip
\texttt{Cliente: Empresa SAC} \\
\texttt{Total: S/ 1,500.00} \\
\texttt{Documento: F001-00001234}

\medskip
\texttt{Te enviamos el comprobante en el siguiente mensaje.}

\medskip
\textit{[Se env\'ia el PDF del comprobante]}
\end{mensajebot}

% ============================================================
% REGISTRO DE COMPRA
% ============================================================
\section{Registro de una compra}

El flujo de compra es similar al de venta con estas diferencias:

\subsection{Diferencias con la venta}

\begin{center}
\begin{tabular}{>{\bfseries}l p{5cm} p{5cm}}
\toprule
Aspecto & Venta & Compra \\
\midrule
Entidad & Cliente & Proveedor \\
Centro de costo & No se pide & S\'i se pide (paso extra) \\
Al finalizar & Texto + PDF & Solo texto \\
Tipo documento & Factura / Boleta / Nota de venta & Factura / Boleta / Nota de compra \\
\bottomrule
\end{tabular}
\end{center}

\subsection{Opciones operativas en compra}

En compras, las opciones son tres (una m\'as que en venta):

\begin{enumerate}
    \item \textbf{Sucursal}
    \item \textbf{Centro de costo} (solo en compras)
    \item \textbf{Forma de pago}
\end{enumerate}

\subsection{Ejemplo de compra}

\begin{mensajeusuario}
Compra de factura por 500 soles al proveedor Juan P\'erez, DNI 12345678, contado
\end{mensajeusuario}

\begin{mensajebot}
\texttt{COMPRA} \\
\texttt{Factura} \\
\texttt{Proveedor: Juan P\'erez} \\
\texttt{DNI: 12345678} \\
\texttt{Total: S/ 500.00} \\
\texttt{Moneda: PEN} \\
\texttt{M\'etodo: Contado}

\medskip
\texttt{Confirmar todo para continuar?}
\end{mensajebot}

% ============================================================
% NOTAS DE VENTA Y COMPRA
% ============================================================
\section{Notas de venta y notas de compra}

Las notas son documentos internos con reglas especiales:

\begin{itemize}
    \item \textbf{No se calcula IGV}: el monto total es el dato principal, sin desglose de impuestos
    \item \textbf{No se pide n\'umero de serie}: el sistema no pregunta por serie-n\'umero del comprobante
    \item \textbf{Documento del cliente/proveedor}: si el monto es menor a S/ 700, el RUC/DNI es \textbf{opcional}
\end{itemize}

\begin{mensajeusuario}
Nota de venta por 200 soles
\end{mensajeusuario}

\begin{mensajebot}
\texttt{VENTA} \\
\texttt{Nota de venta} \\
\texttt{Total: S/ 200.00}

\medskip
\texttt{Por favor, br\'indame estos datos:} \\
\texttt{1. Cliente: nombre o raz\'on social} \\
\texttt{2. Moneda: soles (PEN) o d\'olares (USD)?} \\
\texttt{3. M\'etodo de pago: contado o cr\'edito?}
\end{mensajebot}

Note que \textbf{no pregunta} por RUC/DNI (monto < 700) ni por n\'umero de serie.

% ============================================================
% MÉTODO DE PAGO
% ============================================================
\section{M\'etodo de pago: contado vs cr\'edito}

\subsection{Contado}

Si el usuario indica \texttt{contado}, el bot lo registra y \textbf{no hace preguntas adicionales} sobre d\'ias de cr\'edito ni cuotas.

\subsection{Cr\'edito}

Si el usuario indica \texttt{cr\'edito}, el bot preguntar\'a:

\begin{mensajebot}
\texttt{1. A cu\'antos d\'ias ser\'a el plazo? Puede ser 15, 30, 45, 60 o 90 d\'ias.} \\
\texttt{2. En cu\'antas cuotas ser\'a el pago? (de 1 a 24)}
\end{mensajebot}

\begin{notaimportante}
Una vez registrado el m\'etodo de pago, el bot \textbf{no volver\'a a preguntar} por este campo. Si necesita cambiarlo, simplemente escriba el nuevo valor (por ejemplo: \texttt{"cambiar a cr\'edito"}).
\end{notaimportante}

% ============================================================
% DOCUMENTO ADJUNTO
% ============================================================
\section{Documento adjunto}

Cuando el usuario env\'ia un \textbf{documento PDF} o una \textbf{imagen} por WhatsApp, el bot lo almacena autom\'aticamente como documento adjunto de la operaci\'on:

\begin{itemize}
    \item El archivo se guarda y \textbf{se preserva} durante todo el flujo, sin importar cu\'antos mensajes se env\'ien despu\'es
    \item Se vincula autom\'aticamente al comprobante al finalizar la operaci\'on
    \item Solo se elimina cuando se borra o cancela el registro completo
    \item Si se env\'ia un nuevo documento, reemplaza al anterior
\end{itemize}

\begin{notaimportante}
No es necesario enviar el documento en un momento espec\'ifico del flujo. Se puede enviar al inicio, durante la carga de datos o justo antes de confirmar.
\end{notaimportante}

% ============================================================
% COMANDOS ÚTILES
% ============================================================
\section{Comandos y acciones disponibles}

\begin{center}
\begin{tabular}{>{\bfseries}l p{8cm}}
\toprule
Acci\'on & Qu\'e decir \\
\midrule
Ver resumen & ``Qu\'e llevo?'', ``Dame el resumen'', ``C\'omo va mi registro?'' \\
Confirmar & ``S\'i'', ``Confirmo'', ``Dale'', ``Ok'', ``Listo'' \\
Cancelar & ``Cancelar'', ``Borrar'', ``Empezar de cero'' \\
Cambiar dato & Simplemente escribir el nuevo valor \\
\bottomrule
\end{tabular}
\end{center}

% ============================================================
% ERRORES Y REINTENTOS
% ============================================================
\section{Errores y reintentos}

\subsection{Si el bot no responde}

Si env\'ias un mensaje y \textbf{no recibes ninguna respuesta} (ni de \'exito ni de error), esto indica un \textbf{problema de conexi\'on}, no un fallo del bot. Simplemente \textbf{vuelve a enviar tu mensaje} y el bot continuar\'a el flujo normalmente desde donde qued\'o.

\subsection{Errores al emitir}

Si ocurre un error al emitir el comprobante (por ejemplo, un error de SUNAT o datos inv\'alidos), el bot:

\begin{enumerate}
    \item Muestra el mensaje de error espec\'ifico
    \item Regresa al paso de edici\'on (antes de la confirmaci\'on)
    \item Permite corregir los datos problem\'aticos
    \item El usuario puede volver a confirmar e intentar nuevamente
\end{enumerate}

\begin{mensajebot}
\texttt{Error SUNAT: La fecha de emisi\'on no puede ser anterior a 3 d\'ias.}

\medskip
\texttt{Reinici\'e el flujo al paso de edici\'on. Puedes actualizar los datos y volver a confirmar para intentar nuevamente.}
\end{mensajebot}

\begin{mensajeusuario}
La fecha de emisi\'on es hoy
\end{mensajeusuario}

% ============================================================
% PREGUNTAS FRECUENTES
% ============================================================
\section{Preguntas frecuentes}

\begin{description}[style=nextline]
    \item[Puedo enviar todos los datos en un solo mensaje?]
    S\'i. El bot extrae autom\'aticamente todos los datos que encuentre en el mensaje y solo pregunta por lo que falta.

    \item[Puedo cambiar un dato despu\'es de haberlo ingresado?]
    S\'i. Simplemente escribe el nuevo valor (por ejemplo: ``el monto es 2000'' o ``cambiar moneda a d\'olares'').

    \item[Qu\'e pasa si el bot no responde?]
    Si no recibes un mensaje de error, probablemente se trat\'o de un \textbf{problema de conexi\'on}, no de un fallo del bot. Simplemente vuelve a enviar tu mensaje y el bot continuar\'a desde donde qued\'o.

    \item[Qu\'e pasa si el RUC/DNI no est\'a registrado?]
    Si env\'ias un RUC o DNI junto con el nombre de la persona o empresa, el bot lo \textbf{registrar\'a autom\'aticamente} como cliente (en ventas) o proveedor (en compras) para que el flujo sea m\'as r\'apido en futuras operaciones.

    Si el RUC o DNI ya est\'a registrado, el bot devolver\'a autom\'aticamente el nombre asociado sin necesidad de que lo escribas.

    \item[Qu\'e pasa si hay un error en el nombre del cliente o proveedor?]
    Si hubo un error tipogr\'afico al registrar el nombre, este debe corregirse desde la \textbf{p\'agina web de Maravia} editando directamente al cliente o proveedor. El bot no permite editar entidades ya registradas.

    \item[D\'onde queda registrada la operaci\'on?]
    La venta queda registrada en la \textbf{tabla de ventas} y la compra en la \textbf{tabla de compras} del sistema Maravia, disponibles para consulta desde la plataforma web.

    \item[Puedo registrar varias operaciones a la vez?]
    No. Cada conversaci\'on maneja una operaci\'on a la vez. Debes finalizar o cancelar la actual antes de iniciar otra.

    \item[C\'omo cancelo una operaci\'on?]
    Escribe ``cancelar'', ``borrar'' o ``empezar de cero''. El registro se eliminar\'a por completo.
\end{description}

% ============================================================
% REGISTRO EN EL SISTEMA
% ============================================================
\section{D\'onde queda registrada la operaci\'on}

Al finalizar exitosamente:

\begin{itemize}
    \item Las \textbf{ventas} quedan registradas en la \textbf{tabla de ventas} del sistema Maravia, con su comprobante electr\'onico emitido ante SUNAT (cuando corresponda).
    \item Las \textbf{compras} quedan registradas en la \textbf{tabla de compras} del sistema Maravia.
\end{itemize}

Ambas son consultables desde la \textbf{plataforma web de Maravia}. El registro temporal en el bot se elimina autom\'aticamente una vez que la operaci\'on se completa con \'exito.

% ============================================================
% DIAGRAMA DE FLUJO
% ============================================================
\section{Diagrama de flujo resumido}

\begin{center}
\begin{tcolorbox}[colback=gris, colframe=maravia, width=12cm]
\begin{verbatim}
    Usuario saluda
         |
    [Botones: Compra / Venta]
         |
    Proporciona datos
    (monto, cliente, documento, moneda, metodo)
         |
    Bot pregunta por lo que falta
    (repite hasta completar)
         |
    "Confirmar todo para continuar?"
         |
    Usuario confirma ---------> [OPCIONES]
                                    |
                              Sucursal
                              Centro costo (compra)
                              Forma de pago
                                    |
                        "Brindame una confirmacion"
                                    |
                        Usuario confirma
                                    |
                            [EMISION]
                           /         \
                      VENTA         COMPRA
                   Texto+PDF      Solo texto
                        \         /
                     Registro eliminado
                     (operacion completa)
\end{verbatim}
\end{tcolorbox}
\end{center}

\end{document}
